V prvej fáze projektu boli natrénované a vyhodnotené modely M1 (ResNet50)
a M2 (ResNet50 + CBAM) na 15-triednej klasifikácii jedál, každý v troch
konfiguráciách a s troma nezávislými behmi.

V základnej konfigurácii (zamrznutý backbone, augmentácia) dosiahol M1
testovaciu accuracy \textbf{86,13~\%} a M2 \textbf{89,07~\%}. Tento rozdiel
je štatisticky významný (Welchov t-test, $p = 0{,}020$). Attention vrstva
CBAM umožnila modelu prispôsobiť zamrznuté ImageNet príznaky doméne jedál
a zlepšila najmä triedy, v ktorých M1 zlyhával (pizza, beef\_tartare,
pancakes). Je to však aj za cenu 24-krát viac trénovateľných parametrov,
približne dvojnásobného času trénovania a väčšieho rozptylu medzi behmi.

Oba doplnkové experimenty potvrdili očakávaný smer vplyvu pri oboch
modeloch:

\begin{itemize}
    \item \textbf{augmentácia} zvýšila accuracy o 1,8~p.\,b. (M1) a
          3,5~p.\,b. (M2) -- model s väčšou kapacitou z nej profituje viac,
    \item \textbf{postupné rozmrazovanie vrstiev} zvýšilo accuracy
          o 6,5~p.\,b. (M1) a 2,8~p.\,b. (M2) a bolo najúčinnejším
          opatrením v celom projekte.
\end{itemize}

Po postupnom rozmrazení backbonu sa výhoda CBAM stráca: M1 dosiahol
\textbf{92,60~\%} a M2 91,91~\% a rozdiel nie je významný ($p = 0{,}37$).
Najlepším modelom tejto fázy je preto \textbf{M1 s postupným
rozmrazovaním} -- je najpresnejší, najstabilnejší (std 0,19~p.\,b.)
a trénuje sa výrazne rýchlejšie ako M2. Záver pre praktické použitie je, že
attention mechanizmus CBAM má zmysel hlavne vtedy, keď sa backbone nedá
alebo nechce doladiť (napr. pre obmedzené výpočtové prostriedky). Pri plnom
doladení siete je jednoduchší model rovnako dobrý.

GradCAM vizualizácia pri oboch modeloch ukázala, že sa sieť pri správnych
predikciách sústreďuje na charakteristické časti jedla. Chyby vznikajú
najmä pri obrázkoch, ktoré naozaj obsahujú prvky inej triedy (hranolky na
hot dogu, poleva na šiškách) alebo pri vizuálne príbuzných jedlách
(lasagna a pizza).

Ďalšie kroky projektu zahŕňajú natrénovanie modelov M3 (ConvNeXt),
M4 (vizuálny transformer) a M5 (hybridný model CoAtNet) pri rovnakých
parametroch trénovania a ich porovnanie s M1 a M2.
